\textbf{Phần 1. Trắc nghiệm lựa chọn}
\setcounterpageref{ex}{0}

\begin{ex} %[4D4V19]
$1$ phút bằng bao nhiêu giây?
\choice
{$10$ giây}
{$100$ giây}
{$360$ giây}
{\True $60$ giây}
\loigiai{
$1$ phút $= 60$ giây.\\
\textbf{Đáp án của câu hỏi là D}
}
\end{ex}

\begin{ex} %[4D4V19]
$1$ thế kỉ bằng bao nhiêu năm?
\choice
{$10$ năm}
{\True $100$ năm}
{$1000$ năm}
{$50$ năm}
\loigiai{
$1$ thế kỉ $= 100$ năm.\\
\textbf{Đáp án của câu hỏi là B}
}
\end{ex}

\begin{ex} %[4D4V19]
Năm $1975$ thuộc thế kỉ nào?
\choice
{Thế kỉ XVIII}
{Thế kỉ XIX}
{\True Thế kỉ XX}
{Thế kỉ XXI}
\loigiai{
Năm $1975$ thuộc thế kỉ XX (từ 1901 đến 2000).\\
\textbf{Đáp án của câu hỏi là C}
}
\end{ex}

\begin{ex} %[4D4V19]
Năm $2000$ thuộc thế kỉ nào?
\choice
{\True Thế kỉ XX}
{Thế kỉ XIX}
{Thế kỉ XXI}
{Thế kỉ XXII}
\loigiai{
Năm $2000$ là năm cuối cùng của thế kỉ XX.\\
\textbf{Đáp án của câu hỏi là A}
}
\end{ex}

\begin{ex} %[4D4V19]
$2$ phút $15$ giây bằng bao nhiêu giây?
\choice
{$215$ giây}
{$125$ giây}
{$145$ giây}
{\True $135$ giây}
\loigiai{
$2 \times 60 + 15 = 135$ giây.\\
\textbf{Đáp án của câu hỏi là D}
}
\end{ex}

\begin{ex} %[4D4V19]
$\dfrac{1}{4}$ thế kỉ bằng bao nhiêu năm?
\choice
{$20$ năm}
{\True $25$ năm}
{$30$ năm}
{$50$ năm}
\loigiai{
$100 : 4 = 25$ năm.\\
\textbf{Đáp án của câu hỏi là B}
}
\end{ex}

\begin{ex} %[4D4V19]
Chiến thắng Điện Biên Phủ diễn ra vào năm $1954$. Năm đó thuộc thế kỉ nào?
\choice
{Thế kỉ XIX}
{Thế kỉ XXI}
{\True Thế kỉ XX}
{Thế kỉ XVIII}
\loigiai{
Năm $1954$ thuộc thế kỉ XX.\\
\textbf{Đáp án của câu hỏi là C}
}
\end{ex}

\begin{ex} %[4D4V19]
Một vận động viên chạy cự li $100\text{ m}$ hết $11$ giây, vận động viên thứ hai chạy hết $\dfrac{1}{5}$ phút. Vận động viên nào chạy nhanh hơn?
\choice
{\True Vận động viên thứ nhất}
{Vận động viên thứ hai}
{Hai vận động viên chạy bằng nhau}
{Không so sánh được}
\loigiai{
VĐV 2: $\dfrac{1}{5} \times 60 = 12$ giây.\\
Vì $11$ giây $< 12$ giây nên VĐV 1 chạy nhanh hơn.\\
\textbf{Đáp án của câu hỏi là A}
}
\end{ex}

\begin{ex} %[4D4V19]
Bác Hồ ra đi tìm đường cứu nước vào năm $1911$. Tính đến năm $2026$ đã được bao nhiêu năm?
\choice
{$114$ năm}
{$116$ năm}
{$120$ năm}
{\True $115$ năm}
\loigiai{
$2026 - 1911 = 115$ năm.\\
\textbf{Đáp án của câu hỏi là D}
}
\end{ex}

\begin{ex} %[4D4V19]
Thế kỉ XXI bắt đầu từ năm nào đến năm nào?
\choice
{Từ năm 2000 đến 2099}
{\True Từ năm 2001 đến 2100}
{Từ năm 2000 đến 2100}
{Từ năm 2001 đến 2099}
\loigiai{
Thế kỉ XXI bắt đầu từ năm 2001 đến hết năm 2100.\\
\textbf{Đáp án của câu hỏi là B}
}
\end{ex}

\begin{ex} %[4D4V19]
$3$ thế kỉ $20$ năm bằng bao nhiêu năm?
\choice
{$3200$ năm}
{$3020$ năm}
{\True $320$ năm}
{$302$ năm}
\loigiai{
$3 \times 100 + 20 = 320$ năm.\\
\textbf{Đáp án của câu hỏi là C}
}
\end{ex}

\begin{ex} %[4D4V19]
Vua Quang Trung đại phá quân Thanh vào năm Kỷ Dậu $1789$. Năm đó thuộc thế kỉ nào?
\choice
{\True Thế kỉ XVIII}
{Thế kỉ XVII}
{Thế kỉ XIX}
{Thế kỉ XVI}
\loigiai{
Năm $1789$ thuộc thế kỉ XVIII.\\
\textbf{Đáp án của câu hỏi là A}
}
\end{ex}

\textbf{Phần 2. Trắc nghiệm đúng sai}
\setcounterpageref{ex}{0}

\begin{ex} %[4D4V19]
Xét tính đúng sai về thế kỉ và năm:
\choiceTF
{\True Năm 1010 Lý Thái Tổ dời đô về Thăng Long thuộc thế kỉ XI.}
{\True Năm 938 Ngô Quyền đánh tan quân Nam Hán trên sông Bạch Đằng thuộc thế kỉ X.}
{Năm 2024 thuộc thế kỉ XX.}
{\True Năm 2024 thuộc thế kỉ XXI.}
\loigiai{
\begin{itemize}
    \item Năm 1010: Thế kỉ XI. Đúng.
    \item Năm 938: Thế kỉ X. Đúng.
    \item Năm 2024: Thế kỉ XXI $\neq$ XX. Ý 3 Sai, Ý 4 Đúng.
\end{itemize}
\begin{itemize}
    \item a) \textbf{Đúng}.
    \item b) \textbf{Đúng}.
    \item c) \textbf{Sai}.
    \item d) \textbf{Đúng}.
\end{itemize}
}
\end{ex}

\begin{ex} %[4D4V19]
Xét tính đúng sai về đơn vị thời gian:
\choiceTF
{\True $1$ giờ $= 3600$ giây.}
{\True $2$ phút $30$ giây $= 150$ giây.}
{$\dfrac{1}{2}$ thế kỉ $= 25$ năm.}
{\True $\dfrac{1}{2}$ thế kỉ $= 50$ năm.}
\loigiai{
\begin{itemize}
    \item $1$ giờ $= 60 \times 60 = 3600$ giây. Đúng.
    \item $2 \times 60 + 30 = 150$ giây. Đúng.
    \item $\dfrac{1}{2}$ thế kỉ $= 50$ năm $\neq 25$ năm. Ý 3 Sai, Ý 4 Đúng.
\end{itemize}
\begin{itemize}
    \item a) \textbf{Đúng}.
    \item b) \textbf{Đúng}.
    \item c) \textbf{Sai}.
    \item d) \textbf{Đúng}.
\end{itemize}
}
\end{ex}

\textbf{Phần 3. Trả lời ngắn}
\setcounterpageref{ex}{0}

\begin{ex} %[4D4V19]
Đổi $3$ phút $20$ giây ra giây.
\shortans{200}
\loigiai{
\textbf{Đáp số:} $200$\\
$3 \times 60 + 20 = 200$ giây.
}
\end{ex}

\begin{ex} %[4D4V19]
Đổi $240$ giây ra phút.
\shortans{4}
\loigiai{
\textbf{Đáp số:} $4$\\
$240 : 60 = 4$ phút.
}
\end{ex}

\begin{ex} %[4D4V19]
Năm $1890$ là năm sinh của Chủ tịch Hồ Chí Minh. Năm đó thuộc thế kỉ thứ mấy (ghi bằng số Ả Rập)?
\shortans{19}
\loigiai{
\textbf{Đáp số:} $19$\\
Thế kỉ 19 (XIX).
}
\end{ex}

\begin{ex} %[4D4V19]
Đổi $4$ thế kỉ ra năm.
\shortans{400}
\loigiai{
\textbf{Đáp số:} $400$\\
$4 \times 100 = 400$ năm.
}
\end{ex}

\begin{ex} %[4D4V19]
Một người đi bộ một vòng quanh công viên hết $8$ phút $45$ giây. Tính thời gian đi bộ theo giây.
\shortans{525}
\loigiai{
\textbf{Đáp số:} $525$\\
$8 \times 60 + 45 = 525$ giây.
}
\end{ex}

\begin{ex} %[4D4V19]
Thế kỉ thứ X có bao nhiêu năm nhuận (biết cứ 4 năm lại có 1 năm nhuận)?
\shortans{25}
\loigiai{
\textbf{Đáp số:} $25$\\
$100 : 4 = 25$ năm nhuận.
}
\end{ex}
